\documentclass[10pt,a4paper]{amsart}
\usepackage{amsmath,amssymb,amsthm,graphicx}
\usepackage{verbatim}
\newenvironment{reference}{\comment}{\endcomment}
%\newenvironment{reference}{}{}
\newenvironment{slogan}{\comment}{\endcomment}
\newenvironment{history}{\comment}{\endcomment}

% For commutative diagrams we use Xy-pic
\usepackage[all]{xy}

% We use 2cell for 2-commutative diagrams.
\xyoption{2cell}
% Theorem environments.
%
\theoremstyle{plain}
\newtheorem{theorem}[subsection]{Theorem}
\newtheorem{proposition}[subsection]{Proposition}
\newtheorem{lemma}[subsection]{Lemma}

\theoremstyle{definition}
\newtheorem{definition}[subsection]{Definition}
\newtheorem{example}[subsection]{Example}
\newtheorem{exercise}[subsection]{Exercise}
\newtheorem{situation}[subsection]{Situation}

\theoremstyle{remark}
\newtheorem{remark}[subsection]{Remark}
\newtheorem{remarks}[subsection]{Remarks}

\numberwithin{equation}{subsection}

% Macros
%
\def\lim{\mathop{\mathrm{lim}}\nolimits}
\def\colim{\mathop{\mathrm{colim}}\nolimits}
\def\Spec{\mathop{\mathrm{Spec}}}
\def\Hom{\mathop{\mathrm{Hom}}\nolimits}
\def\Ext{\mathop{\mathrm{Ext}}\nolimits}
\def\SheafHom{\mathop{\mathcal{H}\!\mathit{om}}\nolimits}
\def\SheafExt{\mathop{\mathcal{E}\!\mathit{xt}}\nolimits}
\def\Sch{\mathit{Sch}}
\def\Mor{\mathop{\mathrm{Mor}}\nolimits}
\def\Ob{\mathop{\mathrm{Ob}}\nolimits}
\def\Sh{\mathop{\mathit{Sh}}\nolimits}
\def\NL{\mathop{N\!L}\nolimits}
\def\CH{\mathop{\mathrm{CH}}\nolimits}
\def\proetale{{pro\text{-}\acute{e}tale}}
\def\etale{{\acute{e}tale}}
\def\QCoh{\mathit{QCoh}}
\def\Ker{\mathop{\mathrm{Ker}}}
\def\Im{\mathop{\mathrm{Im}}}
\def\Coker{\mathop{\mathrm{Coker}}}
\def\Coim{\mathop{\mathrm{Coim}}}

% Boxtimes
%
\DeclareMathSymbol{\boxtimes}{\mathbin}{AMSa}{"02}

%
% Macros for moduli stacks/spaces
%
\def\QCohstack{\mathcal{QC}\!\mathit{oh}}
\def\Cohstack{\mathcal{C}\!\mathit{oh}}
\def\Spacesstack{\mathcal{S}\!\mathit{paces}}
\def\Quotfunctor{\mathrm{Quot}}
\def\Hilbfunctor{\mathrm{Hilb}}
\def\Curvesstack{\mathcal{C}\!\mathit{urves}}
\def\Polarizedstack{\mathcal{P}\!\mathit{olarized}}
\def\Complexesstack{\mathcal{C}\!\mathit{omplexes}}
% \Pic is the operator that assigns to X its picard group, usage \Pic(X)
% \Picardstack_{X/B} denotes the Picard stack of X over B
% \Picardfunctor_{X/B} denotes the Picard functor of X over B
\def\Pic{\mathop{\mathrm{Pic}}\nolimits}
\def\Picardstack{\mathcal{P}\!\mathit{ic}}
\def\Picardfunctor{\mathrm{Pic}}
\def\Deformationcategory{\mathcal{D}\!\mathit{ef}}
\usepackage{cleveref}
\crefname{originallemma}{original lemma}{original lemmas}
\Crefname{originallemma}{Original lemma}{Original lemmas}
\makeatletter
\newlabel{lemma-relative-dimension-is-semi-continuous}{{0DRQ}{source}{}{sourcepointer-1}{}}
\newlabel{lemma-relative-dimension-is-semi-continuous@cref}{{[originallemma][0][]0DRQ}{[1][0][]source}{}{}{}}
\makeatother
\crefname{originalreference}{original reference}{original references}
\Crefname{originalreference}{Original reference}{Original references}
\makeatletter
\newlabel{stacks-morphisms-lemma-quasi-compact-permanence}{{050Y}{source}{}{sourcepointer-2}{}}
\newlabel{stacks-morphisms-lemma-quasi-compact-permanence@cref}{{[originalreference][0][]050Y}{[1][0][]source}{}{}{}}
\makeatother
\crefname{originalreference}{original reference}{original references}
\Crefname{originalreference}{Original reference}{Original references}
\makeatletter
\newlabel{stacks-morphisms-lemma-surjection-from-quasi-compact}{{050X}{source}{}{sourcepointer-3}{}}
\newlabel{stacks-morphisms-lemma-surjection-from-quasi-compact@cref}{{[originalreference][0][]050X}{[1][0][]source}{}{}{}}
\makeatother
\crefname{originallemma}{original lemma}{original lemmas}
\Crefname{originallemma}{Original lemma}{Original lemmas}
\makeatletter
\newlabel{lemma-map-of-components}{{0DR5}{source}{}{sourcepointer-4}{}}
\newlabel{lemma-map-of-components@cref}{{[originallemma][0][]0DR5}{[1][0][]source}{}{}{}}
\makeatother
\crefname{originallemma}{original lemma}{original lemmas}
\Crefname{originallemma}{Original lemma}{Original lemmas}
\makeatletter
\newlabel{lemma-irreducible-implies-equidimensional}{{0DRX}{source}{}{sourcepointer-5}{}}
\newlabel{lemma-irreducible-implies-equidimensional@cref}{{[originallemma][0][]0DRX}{[1][0][]source}{}{}{}}
\makeatother
\crefname{originallemma}{original lemma}{original lemmas}
\Crefname{originallemma}{Original lemma}{Original lemmas}
\makeatletter
\newlabel{lemma-behaviour-of-dimensions-wrt-smooth-morphisms}{{0DRI}{source}{}{sourcepointer-6}{}}
\newlabel{lemma-behaviour-of-dimensions-wrt-smooth-morphisms@cref}{{[originallemma][0][]0DRI}{[1][0][]source}{}{}{}}
\makeatother
\crefname{originalreference}{original reference}{original references}
\Crefname{originalreference}{Original reference}{Original references}
\makeatletter
\newlabel{morphisms-definition-relative-dimension-d}{{02NJ}{source}{}{sourcepointer-7}{}}
\newlabel{morphisms-definition-relative-dimension-d@cref}{{[originalreference][0][]02NJ}{[1][0][]source}{}{}{}}
\makeatother
\crefname{originalreference}{original reference}{original references}
\Crefname{originalreference}{Original reference}{Original references}
\makeatletter
\newlabel{morphisms-lemma-rel-dimension-dimension}{{0AFE}{source}{}{sourcepointer-8}{}}
\newlabel{morphisms-lemma-rel-dimension-dimension@cref}{{[originalreference][0][]0AFE}{[1][0][]source}{}{}{}}
\makeatother
\title{Fibre dimension dominates the dimension drop along a dominant morphism, with equality on a dense open}
\author{The Stacks Project Authors}
\date{stacks-project:stacks-geometry.tex:local-snapshot-20261012; GNU FDL 1.2 or later}
\begin{document}
\maketitle
\noindent\textit{Editorial note.} All delivered mathematical text is the author's own, copied verbatim from the pinned Stacks Project chapter file stacks-geometry.tex; no statement or proof step was rewritten, reordered or removed.\par\medskip
\noindent\textit{Editorial note.} Span recomputation by label pairing (LANE\_BRIEF 10.2): in this chapter the author's proof environment follows the statement environment as a standalone block; the spec covers the complete environment pair (statement span, proof span) as recorded in the pool candidate.\par\medskip
\noindent\textit{Editorial note.} Occupation scan: combined claim\_key + statement\_anchor + generic-line scan over all registration files (active and withdrawn) found 0 hits for stacks-geometry.tex (ge-occupied.json, 2026-10-14); 207 route=geometry / chapter-list field hits in focused-seeds are alias noise (route name and alphabetical chapter enumerations), excluded per the ladder alias-blind-spot caveat.\par\medskip
\ifdefined\hypertarget\hypertarget{sourcepointer-1}{}\fi
\noindent\textit{Source reference.} Original lemma 0DRQ: sources/stacks-chow/tags.txt, source lines 17230--17230. Cross-chapter standard prerequisite cited by the author inside the counted span; it is used as a pointer and its own proof is not part of this item.\par\medskip
\ifdefined\hypertarget\hypertarget{sourcepointer-2}{}\fi
\noindent\textit{Source reference.} Original reference 050Y: sources/stacks-chow/tags.txt, source lines 6516--6516. Cross-chapter standard prerequisite cited by the author inside the counted span; it is used as a pointer and its own proof is not part of this item.\par\medskip
\ifdefined\hypertarget\hypertarget{sourcepointer-3}{}\fi
\noindent\textit{Source reference.} Original reference 050X: sources/stacks-chow/tags.txt, source lines 6515--6515. Cross-chapter standard prerequisite cited by the author inside the counted span; it is used as a pointer and its own proof is not part of this item.\par\medskip
\ifdefined\hypertarget\hypertarget{sourcepointer-4}{}\fi
\noindent\textit{Source reference.} Original lemma 0DR5: sources/stacks-chow/tags.txt, source lines 17210--17210. Cross-chapter standard prerequisite cited by the author inside the counted span; it is used as a pointer and its own proof is not part of this item.\par\medskip
\ifdefined\hypertarget\hypertarget{sourcepointer-5}{}\fi
\noindent\textit{Source reference.} Original lemma 0DRX: sources/stacks-chow/tags.txt, source lines 17237--17237. Cross-chapter standard prerequisite cited by the author inside the counted span; it is used as a pointer and its own proof is not part of this item.\par\medskip
\ifdefined\hypertarget\hypertarget{sourcepointer-6}{}\fi
\noindent\textit{Source reference.} Original lemma 0DRI: sources/stacks-chow/tags.txt, source lines 17223--17223. Cross-chapter standard prerequisite cited by the author inside the counted span; it is used as a pointer and its own proof is not part of this item.\par\medskip
\ifdefined\hypertarget\hypertarget{sourcepointer-7}{}\fi
\noindent\textit{Source reference.} Original reference 02NJ: sources/stacks-chow/tags.txt, source lines 3482--3482. Cross-chapter standard prerequisite cited by the author inside the counted span; it is used as a pointer and its own proof is not part of this item.\par\medskip
\ifdefined\hypertarget\hypertarget{sourcepointer-8}{}\fi
\noindent\textit{Source reference.} Original reference 0AFE: sources/stacks-chow/tags.txt, source lines 13157--13157. Cross-chapter standard prerequisite cited by the author inside the counted span; it is used as a pointer and its own proof is not part of this item.\par\medskip
\begin{lemma}
\label{lemma-dims-of-images}
Let $f: \mathcal{T} \to \mathcal{X}$ be a locally of finite type
morphism of Jacobson, pseudo-catenary, and locally Noetherian
algebraic stacks,
whose source is irreducible and whose target is quasi-separated,
and let $\mathcal{Z} \hookrightarrow \mathcal{X}$ denote the scheme-theoretic
image of $\mathcal{T}$.
Then for all $t \in |T|$,
we have that
$\dim_t( \mathcal{T}_{f(t)}) \geq \dim \mathcal{T}  - \dim \mathcal{Z}$,
and there is a non-empty (equivalently, dense)
open subset of $|\mathcal{T}|$ over which equality holds.
\end{lemma}

\begin{proof}
Replacing $\mathcal{X}$ by $\mathcal{Z}$, we may and do assume that $f$ is
scheme theoretically dominant,
and also that $\mathcal{X}$ is irreducible.
By the upper semi-continuity of fibre dimensions
(Lemma \ref{lemma-relative-dimension-is-semi-continuous} (1)),
it suffices to prove that the equality
$\dim_t( \mathcal{T}_{f(t)}) =\dim \mathcal{T}  - \dim \mathcal{Z}$
holds for $t$ lying in
some non-empty open substack of $\mathcal{T}$.
For this reason, in the argument we are always free
to replace $\mathcal{T}$ by a non-empty open substack.

\medskip\noindent
Let $T' \to \mathcal{T}$ be a smooth surjective morphism whose source
is a scheme, and let $T$ be a non-empty quasi-compact open subset
of $T'$.  Since $\mathcal{Y}$ is quasi-separated, we find
that $T \to  \mathcal{Y}$ is quasi-compact
(by Morphisms of Stacks, Lemma
\ref{stacks-morphisms-lemma-quasi-compact-permanence}, applied to the morphisms
$T \to \mathcal{Y} \to \Spec \mathbf{Z}$).
Thus, if we replace $\mathcal{T}$ by the image of $T$ in $\mathcal{T}$,
then we may assume (appealing to
Morphisms of Stacks, Lemma
\ref{stacks-morphisms-lemma-surjection-from-quasi-compact}
that the morphism $f:\mathcal{T} \to \mathcal{X}$ is quasi-compact.

\medskip\noindent
If we choose a smooth surjection $U \to \mathcal{X}$ with $U$ a scheme,
then Lemma \ref{lemma-map-of-components} ensures that
we may find an irreducible open subset $V$ of $U$ such
that $V \to \mathcal{X}$ is smooth and scheme-theoretically dominant.
Since scheme-theoretic dominance for quasi-compact morphisms
is preserved by flat base-change,
the base-change $\mathcal{T} \times_{\mathcal{X}} V \to V$
of the scheme-theoretically
dominant morphism $f$ is again
scheme-theoretically dominant.   We let $Z$ denote a scheme
admitting a smooth surjection onto this fibre product;
then $Z \to \mathcal{T} \times_{\mathcal{X}} V \to V$
is again scheme-theoretically dominant.
Thus we may find an irreducible
component $C$ of $Z$ which scheme-theoretically
dominates $V$.
Since the composite  $Z \to \mathcal{T}\times_{\mathcal{X}} V \to \mathcal{T}$
is smooth,
and since $\mathcal{T}$ is irreducible,
Lemma \ref{lemma-map-of-components} shows that any irreducible
component of the source has dense image in $|\mathcal{T}|$.
We now replace
$C$ by a non-empty open subset $W$ which is disjoint from every other
irreducible component of $Z$, and
then replace $\mathcal{T}$ and $\mathcal{X}$ by the images of $W$
and $V$
(and apply Lemma \ref{lemma-irreducible-implies-equidimensional}
to see that this
doesn't change the dimension of either $\mathcal{T}$ or $\mathcal{X}$).
If we let $\mathcal{W}$ denote the image of the morphism
$W \to \mathcal{T}\times_{\mathcal{X}} V$,
then $\mathcal{W}$ is open in $\mathcal{T}\times_{\mathcal{X}} V$ (since the
morphism $W \to \mathcal{T}\times_{\mathcal{X}} V$ is smooth),
and is irreducible (being the image of an irreducible
scheme).  Thus we end up with a commutative diagram
$$
\xymatrix{
W \ar[dr] \ar[r]  & \mathcal{W} \ar[r] \ar[d] & V \ar[d] \\
& \mathcal{T} \ar[r] & \mathcal{X}
}
$$
in which $W$ and $V$ are schemes,
the vertical arrows are smooth and surjective,
the diagonal arrows and the left-hand
upper horizontal arroware smooth,
and the induced morphism $\mathcal{W} \to \mathcal{T}\times_{\mathcal{X}} V$ is
an open immersion. Using this diagram, together with the definitions
of the various dimensions involved in
the statement of the lemma, we will reduce our verification
of the lemma to the case of schemes, where it is known.

\medskip\noindent
Fix $w \in |W|$ with image $w' \in |\mathcal{W}|$,
image $t \in |\mathcal{T}|$, image $v$ in $|V|$,
and image $x$ in $|\mathcal{X}|$.
Essentially by definition (using the
fact that $\mathcal{W}$ is open in $\mathcal{T}\times_{\mathcal{X}} V$, and that
the fibre of a base-change is the base-change of the fibre),
we obtain the equalities
$$
\dim_v V_x = \dim_{w'} \mathcal{W}_t
$$
and
$$
\dim_t \mathcal{T}_x = \dim_{w'} \mathcal{W}_v.
$$
By Lemma \ref{lemma-behaviour-of-dimensions-wrt-smooth-morphisms}
(the diagonal arrow and right-hand vertical
arrow in our diagram realise $W$ and $V$ as smooth covers by
schemes of the stacks $\mathcal{T}$ and $\mathcal{X}$), we find that
$$
\dim_t \mathcal{T} = \dim_w W - \dim_w W_t
$$
and
$$
\dim_x \mathcal{X} = \dim_v V - \dim_v V_x.
$$
Combining the equalities, we find that
$$
\dim_t \mathcal{T}_x - \dim_t \mathcal{T} + \dim_x \mathcal{X}
= \dim_{w'} \mathcal{W}_v - \dim_w W + \dim_w W_t + \dim_v V -
\dim_{w'} \mathcal{W}_t
$$
Since $W \to \mathcal{W}$ is a smooth surjection, the same is true
if we base-change over the morphism $\Spec \kappa(v) \to V$
(thinking of $W \to \mathcal{W}$ as a morphism over $V$),
and from this smooth morphism we obtain the first of the following
two equalities
$$
\dim_w W_v - \dim_{w'} \mathcal{W}_v = \dim_w (W_v)_{w'} = \dim_w W_{w'};
$$
the second equality follows via a direct comparison of the
two fibres involved.
Similarly, if we think of $W \to \mathcal{W}$ as a morphism of schemes
over $\mathcal{T}$, and base-change over some representative of the point
$t \in |\mathcal{T}|$, we obtain the equalities
$$
\dim_w W_t - \dim_{w'} \mathcal{W}_t = \dim_w (W_t)_{w'} = \dim_w W_{w'}.
$$
Putting everything together, we find that
$$
\dim_t \mathcal{T}_x - \dim_t \mathcal{T} + \dim_x \mathcal{X}
=  \dim_w W_v - \dim_w W + \dim_v V.
$$
Our goal is to show that the left-hand side of this equality
vanishes for a non-empty open subset
of $t$.  As $w$ varies over a non-empty open subset of $W$,
its image $t \in |\mathcal{T}|$ varies over a non-empty open
subset of $|\mathcal{T}|$ (as $W \to \mathcal{T}$ is smooth).

\medskip\noindent
We are therefore reduced to showing that if $W\to V$ is a
scheme-theoretically dominant morphism of irreducible locally
Noetherian schemes that is locally of finite type,
then there is a non-empty open subset of
points $w\in W$ such that $\dim_w W_v =\dim_w W - \dim_v V$
(where $v$ denotes the image of $w$ in $V$).
This is a standard fact,
whose proof we recall for the convenience of the reader.

\medskip\noindent
We may replace $W$ and $V$ by their underlying reduced subschemes
without altering the validity (or not) of this equation,
and thus we may assume that they are in fact integral schemes.
Since $\dim_w W_v$ is locally constant on $W,$ replacing $W$
by a non-empty open subset if necessary, we may assume that $\dim_w W_v$
is constant, say equal to $d$.  Choosing this open subset to be affine,
we may also assume that the morphism $W\to V$ is in fact of finite type.
Replacing $V$ by a non-empty open subset if necessary
(and then pulling back $W$ over this open subset; the resulting pull-back
is non-empty, since the flat base-change of a quasi-compact
and scheme-theoretically
dominant morphism remains scheme-theoretically dominant),
we may furthermore assume that $W$ is flat over $V$.
The morphism $W\to V$ is thus of relative dimension $d$
in the sense of
Morphisms, Definition
\ref{morphisms-definition-relative-dimension-d}
and it follows from
Morphisms, Lemma \ref{morphisms-lemma-rel-dimension-dimension}
that $\dim_w(W) = \dim_v(V) + d,$ as required.
\end{proof}

\end{document}
